% !TEX root = ../main.tex

\section{Endoodwracalny kwantowy silnik Otta}
Kwantowy cykl Otta jest serią naprzemiennych przemian o zerowej pracy i zerowym cieple. Zgodnie z rysunkiem \ref{fig:pv_silnik_otto} zmiany energii po wycałkowaniu mają postać:
\begin{align*}
	Q_{1-2} &= 0, &
	W_{1-2} &= U\pab{T_2,\omega_2} - U\pab{T_1,\omega_1}, \\
	Q_{2-3} &= U\pab{T_3,\omega_2} - U\pab{T_2,\omega_2}, &
	W_{2-3} &= 0, \\
	Q_{3-4} &= 0, &
	W_{3-4} &= U\pab{T_4,\omega_1} - U\pab{T_3,\omega_2}, \\
	Q_{4-1} &= U\pab{T_1,\omega_1} - U\pab{T_4,\omega_2}, &
	W_{4-1} &= 0.
\end{align*}
Parametrem kontrolowanym przez nas jest \(\omega\). Rozpatrujemy przypadek w granicy niskiego oddziaływania w otoczeniem. W tym przypadku stan stacjonarny jest stanem równowagi termodynamicznej opisywanym rozkładem Boltzmanna \cite{endo_otto}. Wzory opisujące \(U\) i \(\vu{rho}\) w bazie wektorów własnych oscylatora harmonicznego maja postać:
\begin{gather*}
	E_n = \hbar\omega\pab{n + \inv{2}}, \\
	Z = \sum_n \exp\pab{-\frac{E_n}{k_B T}} =
	\exp\pab{-\frac{\hbar\omega}{2 k_B T}} \sum_n \bab{\exp\pab{-\frac{\hbar\omega}{k_B T}}}^n =
	\frac{\exp\pab{-\frac{\hbar\omega}{2 k_B T}}}{1-\exp\pab{-\frac{\hbar\omega}{k_B T}}}, \\
	p_n = \inv{Z} \exp\pab{-\frac{E_n}{k_B T}}, \\
	U = \sum_n p_n E_n = \frac{\hbar\omega}{2} \ctgh\pab{\frac{\hbar\omega}{2k_B T}}, \\
	F = - k_B T \ln Z = k_B T \ln\bab{2\sinh\pab{\frac{\hbar\omega}{2k_B T}}}, \\
	S = \inv{T}\pab{U-F} = \frac{\hbar\omega}{2T} \ctgh\pab{\frac{\hbar\omega}{2k_B t}} - k_B \ln\bab{2\sinh\pab{\frac{\hbar\omega}{2k_B T}}}.
\end{gather*}
Funcja \(S\) jest zależna jedynie od \(\frac{\omega}{T}\) i możemy zapisać:
\begin{equation*}
	S\pab{\frac{\omega}{T}} \overset{def}{=} S\pab{T,\omega}.
\end{equation*}
Jest to funkcja monotoniczna, więc istnieje jest ona też odwracalna. Oznacza to, że jeśli wartości funkcji są równe to również argumenty tej funcji są równe.

Podczas przemiany izentropowej entropia na początku i końcu przemianu jest równa, stąd mamy:
\begin{gather*}
	S\pab{T_1,\omega_1} = S\pab{T_2,\omega_2} \quad\Rightarrow\quad
	S\pab{\frac{\omega_1}{T_1}} = S\pab{\frac{\omega_2}{T_2}} \quad\Rightarrow\quad
	\frac{\omega_1}{T_1} = \frac{\omega_2}{T_2}, \\
	S\pab{T_3,\omega_2} = S\pab{T_4,\omega_1} \quad\Rightarrow\quad
	S\pab{\frac{\omega_2}{T_3}} = S\pab{\frac{\omega_1}{T_4}} \quad\Rightarrow\quad
	\frac{\omega_2}{T_3} = \frac{\omega_1}{T_4}.
\end{gather*}
Prace i ciepła w odpowiednich przemianach są równe:
\begin{gather*}
	W_{1-2} = \frac{\hbar}{2}\pab{\omega_2 - \omega_1} \ctgh\pab{\frac{\hbar\omega_1}{2k_B T_1}}, \\
	W_{3-4} = \frac{\hbar}{2}\pab{\omega_1 - \omega_2} \ctgh\pab{\frac{\hbar\omega_2}{2k_B T_3}}, \\
	Q_{2-3} = \frac{\hbar\omega_2}{2} \ctgh\pab{\frac{\hbar\omega_2}{2k_B T_3}} - \frac{\hbar\omega_2}{2} \ctgh\pab{\frac{\hbar\omega_1}{2k_B T_1}}.
\end{gather*}
Wprowadzając wielkość pomocniczą \(\kappa = \frac{\omega_1}{\omega_2}\) znajdujemy sprawność silnika:
\begin{equation}\label{eq:sprawnosc_kwantowy_otto}
	\eta = - \frac{W_{1-2} + W_{3-4}}{Q_{2-3}} = 1 - \kappa.
\end{equation}
Jest to wzór analogiczny do wzoru \eqref{eq:parametry_otto_eta} w tym, że jest on niezależny od temperatur zbiorników.
